% the conclusions

This project successfully proposed the Multi-Domain Inconsistency Fusion (MDIF)
framework for distinguishing AI-generated and AI-manipulated images from
authentic images using a lightweight hybrid architecture. Combining spatial RGB
features, frequency-domain spectral artifacts and depth-shading geometric
inconsistencies into a unified feature representation for image forgery
detection, led to a display of strong overall performance reaching a 92\%
classification accuracy and a 0.98 AUC score. The framework was capable of
exceptional performance on AI-generated images while maintaining a meaningful
performance on inpainted manipulations, despite its difficulty. The results show
an obvious improvement in generalisation and robustness thanks to the
multi-domain fusion strategy, unlike what relying on a single feature domain
would have achieved. Furthermore, the computational efficiency and ease of
deployment in real-world forensic applications, even in resource-constrained
environments, makes this project currently significant and a necessity. A future
scope of this project would be inculcating a larger domain of inpainted images
in dataset to ensure a better accuracy in that case and also test the project
under compression conditions, while extending cross-generator generalisation.
